% !TeX root = ../main.tex
\section{Discussion}
Two decades of ICWSM show substantial platform turnover without erasing familiar research problems. Community research makes this visible: its topic share changes little across the endpoint periods, while governance cues become much more common. The literature shows why that distinction matters. Sampling still determines whose behavior enters a study; traces still need construct validation; causal claims still depend on comparison design; moderation still embeds judgment; and data access still constrains what researchers can observe. AI changes each of these conditions, especially the provenance of language and exposure.

\paragraph{Following questions within topics.}
Topical continuity can conceal changes in what researchers seek to explain or improve. Retrospectives can connect subjects with research questions, designs, and intended beneficiaries. The weak stability of our phase boundaries favors overlapping interests over sharply separated eras. Historical interviews and archival work could explain changes that publication profiles cannot.

\paragraph{Making methodological comparisons reusable.}
Table~\ref{tab:checklist} turns recurring design questions into reporting items. A later reader should be able to judge whether a finding still applies after a platform, access regime, or content-production process changes. The checklist is a reporting aid; it does not certify that an assumption holds.

\begin{table*}[t]
\small
\centering
\begin{tabular}{@{}p{.20\textwidth}p{.73\textwidth}@{}}
\textbf{Reporting item} & \textbf{Information to provide} \\
\toprule
Target population & Who or what is the study about? Describe observed users, communities, languages, dates, and exclusions; distinguish the sample from the target population. \\
\rowcollight Question and estimand & State the descriptive quantity or causal contrast, unit of analysis, exposure, outcome, and time horizon. Specify whose effect is estimated. \\
Constructs and measures & Explain why each trace or label measures the intended construct. Report validation, disagreements, missingness, and relevant subgroup or period differences. \\
\rowcollight Collection and provenance & Record interfaces, sampling rules, access restrictions, deletion or attrition, and data and model versions. Document content-production processes where known. \\
Assignment and timing & Describe how exposure arises, the comparison group, and the ordering of covariates, treatment, and outcomes. Identify concurrent events and possible spillovers. \\
\rowcollight Adjustment and overlap & Justify pretreatment covariates, including text features. Report balance, overlap, exclusions, and any change in the target population after adjustment. \\
Sensitivity and uncertainty & State assumptions and plausible failures. Use appropriate sensitivity analyses, alternative measures, placebo or negative-control comparisons, and uncertainty estimates. \\
\rowcollight Interpretation and reuse & Separate prediction, association, and causal interpretation. State limits of transfer and provide reusable code, definitions, and access information when possible. \\
\bottomrule
\end{tabular}
\caption{A reporting checklist for studies using social media observations. Items should be applied to the study's question and design, with non-applicable items explained. The checklist is a reporting aid, not a quality score or a substitute for validation.}
\label{tab:checklist}
\end{table*}


\paragraph{Interpreting AI-mediated interaction.}
AI-mediated communication includes systems that modify, augment, or generate messages on a person's behalf \citep{hancock2020}. A post can now reflect a person's intentions, model suggestions, interface defaults, and platform rules. Final text alone may no longer reveal who composed a message or whose attention it represents.

This matters for measures that infer attitudes, distress, or social support from text. Experiments with suggested replies show that assistance can change language and interpersonal evaluations, while suspected assistance can lower how a partner is perceived \citep{hohenstein2023}. Language quality, perceived care, and interaction benefit should be measured separately. Our LLM-related publication counts motivate these questions; they do not measure how much social media communication is AI-mediated.

AI mediation also changes exposure. A person may encounter generated replies, ranking decisions, and moderation actions that change over time. Studies can record whether people composed, edited, selected, or only received generated content, together with system versions and interaction context. Consented process records, participant accounts, and interface experiments can provide evidence that final text cannot recover. These records support questions about delegation, trust, social obligation, and community norms even when human and machine contributions cannot be cleanly separated.

\paragraph{Preserving comparisons across changing systems.}
Platform access, model versions, and content-production processes can change both behavior and observation. Recording validation evidence, disagreement, collection conditions, and system versions lets later researchers assess whether an earlier finding transfers. Cumulative comparisons depend on preserving the conditions under which a claim was supported.
